\documentclass[10pt]{article}
\usepackage[margin=0.75in]{geometry}
\usepackage{amsmath,amssymb,amsthm,microtype}
\newtheorem{theorem}{Theorem}
\title{Arbitrarily Long Finite Strings of Growing Returns}
\author{Collatz proof project}
\date{September 5, 2026}
\begin{document}
\maketitle
\begin{abstract}
For every prescribed number of visits, we construct a positive natural
Collatz seed whose next consecutive returns to two modulo nine all grow.
Each return takes four steps with parity word $1101$. This rules out a
uniform fixed-return-count descent argument. The seed depends on the
requested length; no infinite escaping orbit is constructed.
\end{abstract}
Use the shortcut map $T(n)=n/2$ for even $n$ and $(3n+1)/2$ for odd $n$.
\begin{theorem}
For every $R\geq0$ there is a positive $n\equiv2\pmod9$ such that its
first $R$ positive-time returns to that progression occur at times
$4,8,\ldots,4R$, and every return is larger than the preceding one.
There are exactly $3R$ odd steps through time $4R$.
\end{theorem}
\begin{proof}
If $n=16q+11$, exact cylinder transport gives the four iterates
\[
24q+17,\quad36q+26,\quad18q+13,\quad27q+20.
\]
The four steps have parity word $1101$, and $16T^4(n)=27n+23$.
If also $n\equiv2\pmod9$, then $q\equiv0\pmod9$; the residues are
$8,8,4,2$. Thus this is a first return and strictly increases $n$.

If $16^{R+1}$ divides $11n+23$, reduction modulo sixteen forces
$n\equiv11\pmod{16}$. The four-step identity gives
\[
16\bigl(11T^4(n)+23\bigr)=27(11n+23),
\]
so $16^R$ divides $11T^4(n)+23$. Induction therefore realizes $R$
successive blocks whenever $16^R\mid11n+23$ and $n\equiv2\pmod9$.

Define $n_0=2$ and $n_{R+1}=144n_R+299$. Induction proves
\[
11n_R+23=45\cdot144^R=16^R\cdot45\cdot9^R,
\qquad n_R\equiv2\pmod9.
\]
These positive seeds meet the required conditions. Adding the odd counts
of the blocks proves the final count $3R$.
\end{proof}

\paragraph{Variable branches can replenish divisibility.}
For a branch $F_w(n)=(A_wn+B_w)/D_w$, put
$q_w(n)=(A_w-D_w)n+B_w$. Repeating the branch gives
$q_w(F_w(n))=A_wq_w(n)/D_w$, which decreases its two-adic valuation
by the block length. This alone is not a termination measure when branches
change. Take $w=1101$ and $v=101$, with respective triples
$(A,B,D)=(27,23,16)$ and $(9,7,8)$. For $r\geq1$, let
$n=16\cdot64^r-5$ and $y=F_w(n)$. Lean proves the exact certificate
\[
 n\equiv11\pmod{144},\quad y\equiv65\pmod{72},\quad
 11n+23\equiv32\pmod{64},\quad y+7=27\cdot64^r.
\]
The first two congruences put the actual consecutive returns in branches
$w$ and $v$. Thus $v_2(q_w(n))=5$ but $v_2(q_v(y))=6r$.
Even this fixed branch pair has no uniform upper bound on the valuation
reset. A length-debit argument needs to account for that reset using
information in the seed. This does not exclude every possible amortized
potential. The new certificate is
\texttt{Collatz.unbounded\_reserve\_reset}.

\paragraph{Consequence for the proof program.}
The coefficient through $R$ returns is $(27/16)^R$, not a contraction.
Thus no fixed number of returns guarantees descent uniformly over all
positive seeds in this progression. This is compatible with the earlier
sixteen-step bound for a \emph{contracting} first return: none of these
returns contract. It is also compatible with eventual contraction at a
seed-dependent time. The quantifiers are $\forall R\,\exists n_R$, not
$\exists n\,\forall R$; this construction does not prove divergence.

\raggedright
\paragraph{Formal verification.}
Run \texttt{lake build Collatz.GrowingReturns} from \texttt{lean/}.
Both results use exact arithmetic, with no finite-scan premise or external
axiom. No mathematical priority claim is made.
\end{document}
